% Appendix fragment for the corrected residual-core manuscript.
% Insert after \appendix in the corrected manuscript, not the superseded base PDF.
\section{Lean formalization of the residual-core theory}
\label{app:lean-residual-core}
\providecommand{\LeanDecl}[1]{\texttt{\detokenize{#1}}}

\paragraph{Scope and source identity.}
This appendix concerns exact mathematics on a finite observation set, with
categorical intercept columns and strictly positive real diagonal weights.
The source is the independent \texttt{formal/} library, pinned to Lean 4.19.0
and Mathlib commit \texttt{c44e0c8ee63ca166450922a373c7409c5d26b00b}.
Verification refers to a source commit whose complete library build and
per-theorem dependency audit both succeed. It does not attach automatically
to an earlier manuscript PDF or to a later untested edit.

\paragraph{Existence, feasibility and uniqueness.}
\LeanDecl{wls_residual_exists_unique} and
\LeanDecl{wls_fit_exists_unique} establish existence and uniqueness of the
residual and fitted vector on every subspace of the finite observation space.
The proof constructs the weighted normal-equation map into the dual of the
fitted space. Positive definiteness makes that map injective, and equal finite
dimensions make it surjective. The result does not assume a numerical optimizer.
\LeanDecl{wlsFit_iff_residual} identifies minimization with the conjunction
\[
 y-r\in S,\qquad \sum_i w_i r_i z_i=0\quad\hbox{for every }z\in S.
\]
Orthogonality alone is not used as a definition of a correct projection.
No uniqueness of redundant FE coefficient coordinates is required.

\paragraph{Single-level and elimination lemmas.}
The unique-level zero-residual lemma corresponds to
\LeanDecl{singleton_residual_zero}. Actual categorical columns are defined
by \LeanDecl{feColumn}; legal removals are specified by \LeanDecl{PeelingTrace}.
The trace records singleton incidence on the surviving observations, not an
assumed rank or an assumed projection equality.
\LeanDecl{peelMatrix_det_one} derives invertibility of the actual pivot block
from that incidence condition. \LeanDecl{off_core_mem} and
\LeanDecl{peeling_feasible_iff} give the reconstruction and feasibility identities.
These include a one-removal trace as the leaf-elimination special case.
\LeanDecl{core_loss_transport} proves equality of the weighted loss of a core
residual and its zero extension.

\paragraph{Recursive residual-core theorem.}
Let $M_D^W y$ denote the uniquely defined mathematical residual, and let $J_C$
insert a core vector at its original observation coordinates and put zeros
elsewhere. The principal total statement is
\LeanDecl{categorical_residual_core_total}:
\[
 M_D^W y = J_C M_{D_C}^{W_C}(y|_C).
\]
This statement applies to every legal finite removal trace, including an empty
trace or an early-stopped trace. It no longer assumes that a core minimizer has
already been supplied: existence is established by the WLS results above.
\LeanDecl{categorical_reencoded_core_total} additionally permits recoding that
preserves equality of level labels on the realized core sample.

\paragraph{Corrected normal-equation corollary.}
\LeanDecl{core_inner_transport} states, for every original observation vector $z$,
\[
 \langle J_C r_C,z\rangle_W=\langle r_C,z|_C\rangle_{W_C}.
\]
Taking $z$ to be each original FE column gives the corrected normal-equation
transport in FE-column coordinates. This is an exact identity, independent of
any numerical error model. The feasibility theorem separately guarantees that
the reconstructed fitted vector belongs to the original FE space.

\paragraph{Terminal cores, weights and multiple responses.}
\LeanDecl{terminal_trace_exists} proves existence of a terminal trace.
\LeanDecl{terminal_core_unique} proves that its survivor set is independent
of the legal removal order. \LeanDecl{categorical_terminal_core_projection}
combines terminal existence with the total projection identity.
\LeanDecl{categorical_multi_rhs_total} applies the same trace to every response,
even with different strictly positive weights for the different responses.
This is reuse of incidence topology, not reuse of an old weighted solution.
Repeated FE tuples retain their distinct observation identities.

\paragraph{Boundaries.}
The theorem does not cover zero weights, indefinite weights, non-diagonal GLS
weights or heterogeneous-slope columns without separate hypotheses and proofs.
Floating-point accuracy, stopping tolerances, queue implementations, encoding
programs and complexity bounds are not part of this mathematical formalization.
Those exclusions are not outstanding requirements for this appendix.
The mathematical claims are also independent of historical originality.
